\chapter{La carte de l'Asie-Pacifique}\label{ch:nw:the-asia-pacific-map}

En 2015, le régulateur boursier indien a reçu des plaintes selon lesquelles l'installation de \omterm{def:nw:colocation-products-and-how-they-are-sold:colo}{colocation} de la National Stock Exchange envoyait ses données tick par tick
sur TCP, membre par membre, dans l'ordre où les membres s'étaient connectés, de sorte que «~le premier à se connecter au serveur le moins chargé
obtiendrait un avantage pour recevoir les données plus vite que les autres~». Son ordonnance de 2019 décrivait la conception en détail et
enjoignait à la bourse de restituer 624.89~crores de roupies, intérêts compris~; le tribunal d'appel a annulé la restitution, le régulateur
a fait appel, et le 18 septembre 2026 la Cour suprême a autorisé un règlement d'environ 1\,492~crores de roupies qui a clos l'affaire. Aucun câble
n'était plus court et aucun serveur n'était plus rapide~: l'avantage était inscrit dans l'ordre dans lequel un ordinateur envoyait le même message à plusieurs destinataires.

Ce chapitre situe les plateformes de la zone Asie-Pacifique (Tokyo, Hong Kong, Singapour, Sydney, Mumbai) avec la méthode des chapitres~10 et~11, puis
étudie le cas~: ce que la diffusion séquentielle fait au moment où chaque membre voit un tick, et ce que changent un randomiseur et un second serveur,
et pourquoi le multicast règle le problème. Les distances dans cette région se mesurent en milliers de kilomètres, et plusieurs plateformes ne publient pas
l'adresse de leur bâtiment~: la table des sites porte donc sa précision avec elle.

\section{Tokyo}

JPX exploite sa \omterm{def:nw:colocation-products-and-how-they-are-sold:colo}{colocation} dans ce qu'elle appelle le JPX Primary Site, où les participants de marché, les fournisseurs d'informations et les éditeurs de logiciels indépendants
installent leurs équipements à côté des systèmes de négociation de la Tokyo Stock Exchange et de l'Osaka Exchange. JPX ne publie pas
l'adresse du bâtiment~; son réseau, arrownet, atteint les mêmes systèmes de négociation depuis des points d'accès situés hors du site. La table des sites place donc Tokyo au niveau de la ville, ce qui introduit une erreur de quelques kilomètres (quelques dizaines de microsecondes),
sans conséquence à l'échelle de cette carte~: Tokyo est à \qty{2880}{\kilo\metre} de Hong Kong, soit \qty{9.6}{\milli\second} de lumière dans le vide et
\qty{14.0}{\milli\second} dans une fibre droite.

\begin{dated}{2026-09}{Où correspondent les plateformes de l'Asie-Pacifique}\label{dat:nw:the-asia-pacific-map:sites}
\begin{itemize}
\item \textbf{Tokyo}~: le JPX Primary Site accueille la \omterm{def:nw:colocation-products-and-how-they-are-sold:colo}{colocation} à côté des systèmes de négociation de la TSE et de l'OSE~; le bâtiment n'est pas publié (niveau de la ville).
\item \textbf{Hong Kong}~: le centre de données de la HKEX, situé dans la Tseung Kwan O Industrial Estate (1 Chun Ying Street), accueille sa \omterm{def:nw:colocation-products-and-how-they-are-sold:colo}{colocation}.
\item \textbf{Singapour}~: la SGX déplace son centre de données principal vers un nouveau bâtiment, que les membres peuvent utiliser pour la \omterm{def:nw:colocation-products-and-how-they-are-sold:colo}{colocation}, en même temps que sa nouvelle
  plateforme de négociation, Iris-ST~; son centre de données secondaire ne change pas~; la certification obligatoire se déroule de janvier à avril 2027 (niveau de la ville).
\item \textbf{Sydney}~: l'ASX fait tourner ses marchés primaires et ses plateformes post-marché depuis l'Australian Liquidity Centre, à Gore Hill, depuis
  février 2012.
\item \textbf{Mumbai}~: la \omterm{def:nw:colocation-products-and-how-they-are-sold:colo}{colocation} de la NSE se situe sur son site Exchange Plaza, dans le Bandra Kurla Complex, avec plus de 200 membres en baies et environ 100
  autres par la \omterm{def:nw:colocation-products-and-how-they-are-sold:colo}{colocation} en tant que service.
\end{itemize}
\end{dated}

\section{Hong Kong et Singapour}

La HKEX propose la \omterm{def:nw:colocation-products-and-how-they-are-sold:colo}{colocation}, qu'elle appelle services d'hébergement, dans son propre centre de données de la Tseung Kwan O Industrial Estate, du côté est
de Hong Kong. La SGX est au milieu d'un changement plus vaste. Ses questions-réponses sur Iris-ST, sa nouvelle plateforme de négociation, indiquent que la
plateforme entrera en service dans un centre de données principal nouvellement construit, que les membres pourront aussi utiliser pour la \omterm{def:nw:colocation-products-and-how-they-are-sold:colo}{colocation}, avec le moteur de négociation organisé selon un dispositif «~2+1~» (une instance principale et un secours à chaud dans le centre de données principal, une troisième instance sur le site secondaire).
Elles répondent aussi à une question
que traite ce chapitre~: depuis la \omterm{def:nw:colocation-products-and-how-they-are-sold:colo}{colocation}, les membres pourront recevoir les données de marché ITCH à la fois en TCP et en multicast.

\begin{table}[ht]
\centering\small
\begin{tabular}{llrrr}
\toprule
De & À & Géodésique & Vide & Fibre \\
 & & (km) & (ms) & (ms) \\
\midrule
Tokyo & Hong Kong TKO & 2\,880.3 & 9.61 & 14.05 \\
Hong Kong TKO & Singapour & 2\,577.1 & 8.60 & 12.57 \\
Tokyo & Singapour & 5\,311.0 & 17.72 & 25.90 \\
Singapour & Sydney ALC & 6\,296.4 & 21.00 & 30.71 \\
Tokyo & Sydney ALC & 7\,784.2 & 25.97 & 37.96 \\
Singapour & Mumbai BKC & 3\,903.3 & 13.02 & 19.04 \\
Hong Kong TKO & Mumbai BKC & 4\,317.3 & 14.40 & 21.05 \\
Tokyo & Mumbai BKC & 6\,743.1 & 22.49 & 32.88 \\
\bottomrule
\end{tabular}
\caption{Distances et \omterm{def:nw:the-north-american-map:floor}{planchers de latence} à sens unique entre les sites de l'Asie-Pacifique présentés dans l'\cref{dat:nw:the-asia-pacific-map:sites}. Tokyo et Singapour
sont placés au niveau de la ville, les autres au niveau de la rue, du quartier ou du district~: erreurs de quelques kilomètres, quelques dizaines de microsecondes au plus, face à
des planchers de l'ordre de la milliseconde. Données~: \texttt{nw\_asia.pair\_rows()}.}\label{tab:nw:the-asia-pacific-map:pairs}
\end{table}

\section{Sydney}

L'ASX a transféré sa plateforme de négociation actions dans l'Australian Liquidity Centre, un centre de données situé à Gore Hill, dans le North Shore de Sydney, en
février 2012. Sa page énonce en une phrase les deux règles du chapitre~9~: chaque client reçoit des \omterm{def:nw:colocation-products-and-how-they-are-sold:mmr}{interconnexions} «~de longueur identique~» vers la
plateforme de négociation, où qu'il se trouve dans le bâtiment, et le bâtiment est neutre vis-à-vis des opérateurs, de sorte que ses clients peuvent amener leurs propres réseaux.
Sydney est le centre le plus isolé de la région~: \qty{6296}{\kilo\metre} de Singapour (\qty{21.0}{\milli\second} dans le vide) et
\qty{7784}{\kilo\metre} de Tokyo. Une société de trading installée à Sydney est en concurrence avec d'autres sociétés du même bâtiment pour les marchés de l'ASX, et, pour tout le reste, n'a personne à proximité.
%

\begin{omfigure}
\begin{tikzpicture}
\begin{axis}[axis equal image,width=11.5cm,xmin=-7200,xmax=5200,ymin=-5000,ymax=4300,xtick distance=2000,xlabel={km à l'est de 120\textdegree{} E},
  ylabel={km au nord de 5\textdegree{} N},tick label style={font=\footnotesize,/pgf/number format/1000 sep={\,}},label style={font=\small},grid=major,grid style={gray!20},
  table/col sep=comma,unbounded coords=jump,scaled ticks=false]
\addplot[omDef,thick,no markers] table[x=x,y=y]{figdata/networks/12-the-asia-pacific-map/map_links.csv};
\addplot[only marks,mark=*,mark size=2.4pt,omThm,visualization depends on={value \thisrow{anchor}\as\anc},nodes near coords,
  point meta=explicit symbolic,every node near coord/.append style={font=\footnotesize,anchor=\anc,black}]
  table[x=x,y=y,meta=name,col sep=comma,row sep=\\]{id,x,y,name,anchor\\ tokyo,2191.7,3392.1,Tokyo (JPX),south\\ tko,-635.1,1911.2,Hong Kong TKO (HKEX),south west\\ singapore,-1794.4,-402.8,Singapour (SGX),north\\ alc,3458.2,-4291.9,Sydney ALC (ASX),east\\ bkc,-5228.2,1554.7,Mumbai BKC (NSE),south\\};
\addplot[only marks,mark=none,nodes near coords,point meta=explicit symbolic,
  every node near coord/.append style={font=\scriptsize,fill=white,inner sep=1pt,anchor=center,text=omMeth}]
  table[x=x,y=y,meta=text]{figdata/networks/12-the-asia-pacific-map/map_labels.csv};
\end{axis}
\end{tikzpicture}

\omcaption{Les plateformes de l'Asie-Pacifique à partir des coordonnées calculées, chaque lien étant étiqueté par sa \omterm{def:nw:the-north-american-map:geodesic}{distance géodésique} et son temps de lumière dans le vide. Une
projection équirectangulaire centrée sur 5\textdegree{}~N, 120\textdegree{}~E déforme fortement les formes sur ces distances~: elle ne sert qu'au dessin,
et chaque nombre de la carte est calculé sur l'ellipsoïde. Données~: \texttt{fig\_asia.py} sur \texttt{firm.geomap}.}\label{fig:nw:the-asia-pacific-map:map}
\end{omfigure}

\section{Mumbai et le cas de la colocation}

\begin{definition}[Diffusion unicast]\label{def:nw:the-asia-pacific-map:unicast}
La \emph{diffusion unicast}\index{diffusion unicast} est la distribution de données de marché par l'envoi d'une copie distincte de chaque message à
chaque destinataire, sur une connexion par destinataire (typiquement TCP), de sorte qu'un émetteur transmet les copies l'une après l'autre.
\end{definition}

Le flux tick par tick de la NSE, selon l'ordonnance de la SEBI, était diffusé en TCP depuis juin 2010 pour le segment des futures et des options, et
le multicast n'est arrivé qu'en avril 2014 (novembre 2014 sur le marché au comptant). L'ordonnance décrit la conception à deux niveaux. Un centre de diffusion
envoyait chaque lot de ticks aux serveurs de distribution (POP), un par un, dans l'ordre où ces serveurs s'étaient connectés
ce jour-là. Chaque serveur POP faisait fonctionner trois processus d'émission, appelés ports, et chaque port envoyait à sa propre liste de membres, l'un après l'autre, dans l'ordre
où les membres s'étaient connectés. Un membre connecté le premier, sur un port d'un serveur connecté le premier, était premier dans
les deux files~; un membre placé sur un serveur peu chargé, tel qu'un serveur de secours destiné à la continuité d'activité, avait moins de membres devant lui.

\begin{omfigure}
\begin{tikzpicture}[font=\footnotesize,b/.style={draw,rounded corners=2pt,minimum height=0.55cm,align=center,inner sep=2pt},>=Stealth]
\node[b,fill=omThm!20,minimum height=1.6cm] (c) at (0,0) {centre\\ de diffusion};
\foreach \k/\y/\lab in {1/1.8/{serveur 1\\ (connecté en premier)},2/0.6/{serveur 2},3/-0.6/{serveur 3},4/-1.8/{secondaire\\ (peu chargé)}}{
  \node[b,fill=omDef!15,minimum width=2.2cm] (s\k) at (3.4,\y) {\lab};
  \draw[->] (c.east) -- node[pos=0.55,fill=white,inner sep=1pt,font=\scriptsize]{\k} (s\k.west);}
\foreach \p/\y in {1/2.4,2/1.8,3/1.2}{
  \node[b,fill=gray!10] (p\p) at (6.4,\y) {port \p};
  \draw[->] (s1.east) -- (p\p.west);}
\node[b,fill=white,minimum width=4.3cm,anchor=west] (m1) at (7.4,2.4) {membres dans l'ordre de connexion~: 1\up{er}, 2\up{e}, \dots, $m$-\up{e}};
\draw[->] (p1.east) -- (m1.west);
\node[align=left,font=\scriptsize,anchor=west] at (7.4,1.5) {chaque port envoie le même tick à ses membres\\ l'un après l'autre, dans
  l'ordre où ils se sont connectés};
\node[align=left,font=\scriptsize,anchor=west,omMeth] at (5.4,-0.9) {multicast~: une seule copie, répliquée par les\\ commutateurs vers chaque
  membre à la fois};
\end{tikzpicture}

\omcaption{L'architecture unicast à deux niveaux du flux TCP tick-by-tick de NSE, telle que la décrit l'ordonnance de la SEBI de 2019, de façon schématique~: le centre envoie
à ses serveurs de distribution dans leur ordre de connexion, et les ports de chaque serveur envoient à leurs membres dans l'ordre de connexion des membres. Le multicast remplace ces deux
files par une seule copie que le réseau réplique.}\label{fig:nw:the-asia-pacific-map:arch}
\end{omfigure}

\begin{proposition}[Unicast séquentiel]\label{prop:nw:the-asia-pacific-map:seq}
Si le centre dépense $c_0$ par copie vers un serveur et un port dépense $c$ par copie vers un membre (multiplié par un facteur $\sigma_s$ pour un serveur $s$ plus lent),
le membre de rang de connexion $r$ ($r = 1, 2, \dots$) sur un port du serveur qui occupe le rang $q$ dans l'ordre du centre reçoit un tick après
\[
  t(q, r) = q\,c_0 + r\,c\,\sigma_s + \varepsilon,
\]
où $\varepsilon$ est un petit retard aléatoire. Avec $m$ membres par port, le membre médian attend environ $c\,m/2$ de plus que le premier sur le même
port, quel que soit le réseau.
\end{proposition}

\begin{proof}
Les copies du centre partent les unes après les autres, donc le $q$-ième serveur reçoit le tick après $q\,c_0$~; le port envoie ensuite $r$ copies avant
celle du $r$-ième membre. Le rang médian sur un port de $m$ membres vaut environ $m/2$.
\end{proof}

\begin{definition}[Équité de la diffusion]\label{def:nw:the-asia-pacific-map:fair}
L'\emph{équité de la diffusion}\index{equite de la diffusion@équité de la diffusion} est la propriété d'une distribution de données de marché selon laquelle les destinataires d'un même
service reçoivent chaque message au même moment, à un écart près qui ne dépend pas de leur identité~: elle se mesure par l'écart des temps de réception
d'un même message et par la question de savoir si l'ordre des destinataires se maintient d'un message au suivant.
\end{definition}

La circulaire de la SEBI de mai 2015 a donné la forme réglementaire à cette règle~: les bourses doivent offrir aux membres colocalisés «~un accès équitable et égal aux
installations et aux flux de données~», doivent veiller à ce qu'ils «~subissent une latence similaire par rapport à l'infrastructure fournie par la bourse~»,
et doivent publier des rapports trimestriels sur la latence. La définition comporte deux volets parce que l'affaire en comportait deux. L'écart indique de combien
le premier destinataire est plus rapide que le dernier. La persistance indique si c'est toujours le même~: un ordre aléatoire avec un grand écart donne à chacun
la chance d'être premier~; un ordre fixe réserve à un seul membre toutes les courses.

\omcode{code/firm/dissem/firm_dissem.py}{47}{64}{Diffusion unicast à deux niveaux~: les copies du centre vers les serveurs dans l'ordre des serveurs, les copies de chaque port
vers ses membres dans l'ordre de connexion, et le randomiseur qui tire à nouveau les deux ordres à chaque tick.}\label{lst:nw:the-asia-pacific-map:unicast}

\begin{omfigure}
\begin{tikzpicture}
\begin{axis}[width=11cm,height=6.2cm,xmin=0.5,xmax=10.5,ymin=0,ymax=130,xlabel={rang de connexion du membre sur son port},
  ylabel={délai moyen de réception (\si{\micro\second})},tick label style={font=\footnotesize},label style={font=\small},grid=major,
  grid style={gray!25},legend columns=4,legend style={font=\footnotesize,at={(0.5,-0.25)},anchor=north,draw=none},
  table/col sep=comma,xtick={1,2,3,4,5,6,7,8,9,10}]
\addplot[omDef,thick,mark=*,mark size=1.5pt] table[x=rank,y=s1]{figdata/networks/12-the-asia-pacific-map/rank_curve.csv};
\addplot[omThm,thick,mark=square*,mark size=1.5pt] table[x=rank,y=s2]{figdata/networks/12-the-asia-pacific-map/rank_curve.csv};
\addplot[omMeth,thick,mark=triangle*,mark size=1.8pt] table[x=rank,y=s3]{figdata/networks/12-the-asia-pacific-map/rank_curve.csv};
\addplot[gray,thick,mark=diamond*,mark size=1.8pt] table[x=rank,y=s4]{figdata/networks/12-the-asia-pacific-map/rank_curve.csv};
\legend{serveur 1\up{er} dans l'ordre du centre,2\up{e},3\up{e},4\up{e}}
\end{axis}
\end{tikzpicture}

\omcaption{Simulation, et non mesure d'une plateforme quelconque~: délai moyen de réception d'un tick selon le rang de connexion du membre sur son port, pour les quatre
serveurs chargés du modèle (trois ports de dix membres chacun~; \qty{5}{\micro\second} par copie au centre, \qty{10}{\micro\second} par copie à un port,
gigue exponentielle de moyenne \qty{1}{\micro\second}~; 200 ticks). Le délai de chaque membre est fixé par deux ordres de connexion. Données~:
\texttt{fig\_asia.py}, \texttt{nw\_asia.rank\_curve()}.}\label{fig:nw:the-asia-pacific-map:rank}
\end{omfigure}

Les \Cref{fig:nw:the-asia-pacific-map:rank,tab:nw:the-asia-pacific-map:designs} exécutent le modèle avec des paramètres choisis pour l'illustration, et non tirés
de l'ordonnance. Dans un unicast simple, le premier membre reçoit un tick après \qty{16}{\micro\second} et le membre médian après
\qty{66}{\micro\second}~; l'ordre est le même d'un tick à l'autre (une corrélation de rang de 1.00), si bien que le membre premier sur un tick devance le
membre médian sur chacun des suivants. Un randomiseur, que l'ordonnance dit avoir été mis au point par NSE en 2011 mais non implémenté dans la diffusion de ce flux, tire à nouveau l'ordre
à chaque tick~: l'écart ne change pas, mais la corrélation de rang tombe à 0.07. Le serveur peu chargé continue de favoriser certains membres~: dans le modèle, le
premier membre d'un tick randomisé est en général l'un des six membres du serveur secondaire, et il devance le membre médian dans 88~\% des cas. Fermer le
serveur secondaire au trafic normal ramène ce chiffre à 51~\%, soit pile ou face~; le multicast réduit l'écart lui-même, d'environ \qty{108}{\micro\second} à la
gigue.

\begin{table}[ht]
\centering\small
\begin{tabular}{lrrrrr}
\toprule
Conception & Premier & Médian & Écart & Corrélation & Premier devance \\
 & (\si{\micro\second}) & (\si{\micro\second}) & (\si{\micro\second}) & de rang & le médian \\
\midrule
Unicast, ordre de connexion fixe & 16.0 & 65.6 & 107.7 & 1.00 & 100~\% \\
Unicast, randomisé & 15.1 & 70.7 & 110.2 & 0.07 & 88~\% \\
Unicast, randomisé, sans secondaire & 15.2 & 69.0 & 106.9 & 0.01 & 51~\% \\
Multicast & 10.0 & 10.7 & 8.4 & -- & 50~\% \\
\bottomrule
\end{tabular}
\caption{Simulation de quatre conceptions de diffusion (les paramètres de la \cref{fig:nw:the-asia-pacific-map:rank}~; 126 membres, dont six sur le serveur
secondaire, que la quatrième conception ferme)~: pour un tick, le premier, le médian et l'écart des temps de réception, la corrélation de rang entre
ticks consécutifs, et la fréquence à laquelle le membre premier sur un tick devance le membre médian de ce tick sur les 200 ticks suivants, les deux réagissant en
\qty{20}{\micro\second} avec un écart type de \qty{5}{\micro\second}. Données~: \texttt{nw\_asia.scenarios()}.}\label{tab:nw:the-asia-pacific-map:designs}
\end{table}

\omcode{code/networks/12-the-asia-pacific-map/python/nw_asia.py}{70}{78}{Chaque conception confrontée aux deux volets de l'équité de la diffusion~: l'écart
d'un tick, et la question de savoir si le premier membre reste premier.}\label{lst:nw:the-asia-pacific-map:designs}

\begin{dated}{2026-09}{Le dossier de colocalisation de NSE}\label{dat:nw:the-asia-pacific-map:case}
Plaintes auprès de la SEBI en 2015~; ordonnance de la SEBI du 30 avril 2019 enjoignant à NSE de restituer 624.89~crore de roupies, intérêts à 12~\% par an
à compter du 1\up{er} avril 2014~; décision ultérieure du Securities Appellate Tribunal annulant la restitution, que la SEBI a contestée en appel~; acceptation
par la SEBI, en juillet 2026, des propositions de NSE pour régler les questions de colocalisation et de fibre noire moyennant environ 1\,492~crore de roupies
(environ 1\,224~crore pour la colocalisation)~; et ordonnance de la Cour suprême du 18 septembre 2026 autorisant le règlement et clôturant les appels de la SEBI.
\end{dated}

La leçon pour une société de trading n'est pas seulement historique. Partout où une plateforme ou un fournisseur distribue des données sur une connexion par client (un
flux TCP, un flux WebSocket au chapitre~20, un relais cloud au chapitre~16), l'ordre des copies constitue une latence qu'aucun câble ne peut corriger, et les
questions de cette section s'appliquent~: dans quel ordre les copies sont-elles envoyées, cet ordre change-t-il, et qui se trouve sur le serveur le moins chargé.

\section{Chine continentale et ses restrictions}

Le régulateur des valeurs mobilières chinois traite la colocalisation comme une question d'équité. Ses règles administratives relatives au trading programmé sur le
marché des valeurs mobilières, publiées en mai 2024 et en vigueur depuis le 8 octobre 2024, exigent qu'une bourse qui propose la colocalisation respecte les principes
de «~sécurité, équité et raisonnabilité~» et répartisse équitablement les ressources, et qu'un courtier veille à ce que le trading reste équitable entre ses clients~; la
note d'accompagnement des bourses cible pour surveillance les comptes qui atteignent 300 ordres ou annulations par seconde, ou 20\,000 par jour. Une
consultation d'avril 2023 sur la colocalisation des serveurs de trading proposait que les bourses gèrent la colocalisation dans des zones dédiées et maintiennent
la latence de communication entre ces zones et les processeurs frontaux du système de négociation central à un niveau constant~: la règle des câbles égaux du
chapitre~9, formulée comme une exigence de latence.

\begin{remark}[Ce que la carte laisse de côté]\label{rem:nw:the-asia-pacific-map:leftout}
Le chapitre ne place pas sur la carte les bâtiments des bourses chinoises continentales~: leur colocalisation passe par des courtiers, selon les règles
ci-dessus, et le tableau des sites ne retiendrait qu'une ville. Les chapitres~22 à~26 ajoutent des plateformes par classe d'actifs, avec la même règle~: aucune ligne
sans source, et une note de précision lorsque la source ne donne qu'une ville.
\end{remark}

\begin{method}[Audit du dispositif de diffusion d'une plateforme]\label{met:nw:the-asia-pacific-map:audit}
\begin{enumerate}
\item Vérifier comment le flux est livré en colocalisation~: multicast, unicast (TCP), ou les deux, comme SGX le proposera sur sa nouvelle plateforme.
\item En cas d'unicast, demander dans quel ordre les copies sont envoyées (par identifiant de connexion, par connexion, ou aléatoirement par message), combien de destinataires partagent un même émetteur,
  et si certains serveurs desservent moins de destinataires.
\item Mesurer~: horodater les mêmes messages sur plusieurs connexions et serveurs (chapitre~5), puis calculer l'écart et la corrélation de rang entre
  les messages.
\item Comparer avec les rapports de latence publiés par la plateforme et ses engagements d'équité~; signaler à la plateforme tout ordre persistant.
\end{enumerate}
\end{method}

\section{Tutoriel~: ordre de diffusion}\label{tut:nw:the-asia-pacific-map:tut}

\begin{tutorial}
\textbf{Objectif.} Étendre la table des sites à l'Asie, puis simuler la \omterm{def:nw:the-asia-pacific-map:unicast}{diffusion unicast} et multicast et mesurer leur équité.
\textbf{État final~:} \cref{fig:nw:the-asia-pacific-map:map,fig:nw:the-asia-pacific-map:rank,fig:nw:the-asia-pacific-map:sorted} et
\cref{tab:nw:the-asia-pacific-map:pairs,tab:nw:the-asia-pacific-map:designs}.
\begin{tutsteps}
\item \textbf{Sites.} Cinq lignes dans \texttt{data/networks/sites.csv}, chacune avec sa précision dans son nom~; \texttt{nw\_asia.pair\_rows()} calcule
  le \cref{tab:nw:the-asia-pacific-map:pairs}.
\item \textbf{Topologie.} \texttt{nw\_asia.topology(m)} construit quatre serveurs de trois ports pour $m$ membres et un serveur secondaire de trois ports pour
  deux~; les coûts sont dans \texttt{PARAMS}.
\item \textbf{Un tick.} \texttt{firm\_dissem.unicast\_us} (\cref{lst:nw:the-asia-pacific-map:unicast}) et \texttt{multicast\_us} donnent l'heure de réception de
  chaque membre~; \texttt{fairness} les résume.
\item \textbf{Nombreux ticks.} \texttt{rank\_stability} et \texttt{race\_win} mesurent la persistance et sa valeur~; \texttt{scenarios()}
  (\cref{lst:nw:the-asia-pacific-map:designs}) remplit le \cref{tab:nw:the-asia-pacific-map:designs}.
\end{tutsteps}
\textbf{Pour aller plus loin.} Donner au serveur secondaire la même charge que les autres~; diviser par deux le coût par copie~; placer 20 membres sur chaque port et
vérifier la \cref{prop:nw:the-asia-pacific-map:seq}.
\end{tutorial}

\begin{omfigure}
\begin{tikzpicture}
\begin{axis}[width=11cm,height=6cm,xmin=0,xmax=100,ymin=0,ymax=175,xlabel={membres, triés par heure de réception (\%)},
  ylabel={délai de réception (\si{\micro\second})},tick label style={font=\footnotesize},label style={font=\small},grid=major,
  grid style={gray!25},legend pos=north west,legend style={font=\footnotesize},table/col sep=comma]
\addplot[omDef,thick,no markers] table[x=pct,y=us]{figdata/networks/12-the-asia-pacific-map/sorted_unicast.csv};
\addplot[omMeth,thick,dashed,no markers] table[x=pct,y=us]{figdata/networks/12-the-asia-pacific-map/sorted_randomised.csv};
\addplot[omThm,thick,no markers] table[x=pct,y=us]{figdata/networks/12-the-asia-pacific-map/sorted_multicast.csv};
\legend{unicast,{unicast, aléatoire, sans serveur secondaire},multicast}
\end{axis}
\end{tikzpicture}

\omcaption{Simulation~: délais de réception d'un tick, membres triés du premier au dernier, selon trois conceptions du
\cref{tab:nw:the-asia-pacific-map:designs}. Randomiser change qui se trouve où sur la courbe, pas la courbe~; le multicast l'aplatit. Données~:
\texttt{fig\_asia.py}.}\label{fig:nw:the-asia-pacific-map:sorted}
\end{omfigure}

\section{Construction~: le modèle de diffusion}\label{bld:nw:the-asia-pacific-map:dissem}

\begin{build}
\bfield{Rôle} Un modèle de distribution des données de marché d'une plateforme, pour auditer une conception documentée (chapitre~22) et pour expliquer un écart mesuré
entre connexions.
\bfield{Interface} \texttt{firm\_dissem}~: \texttt{Server(name, ports, speed)}, \texttt{Member(server, port, rank)}, \texttt{members},
\texttt{unicast\_us(servers, c\_centre, c\_send, jitter, rng, shuffle)}, \texttt{multicast\_us}, \texttt{fairness}, \texttt{rank\_stability},
\texttt{race\_win}.
\bfield{Règles} Chaque coût est un paramètre et chaque sortie est étiquetée comme une simulation~; le modèle reproduit l'ordre d'envoi documenté, et non la
latence mesurée d'une plateforme.
\bfield{Tests d'acceptation} \texttt{code/firm/dissem/tests/}~: heures de réception calculées à la main pour une petite topologie, multicast plat, un ordre fixe stable
et un ordre aléatoire non stable, et les probabilités de gagner la course à 0 et à 50 microsecondes.
\bfield{Pour aller plus loin} Les effets propres à TCP (la fenêtre d'un récepteur lent bloquant son port, les retransmissions)~; un commutateur qui réplique le multicast
port par port.
\end{build}

\begin{omsources}
\item SEBI, Order in the matter of NSE colocation, 30 avril 2019~; SEBI circular CIR/MRD/DP/07/2015 on co-location and proximity hosting, 13 mai
2015~; Business Standard, 18 septembre 2026, sur l'ordonnance de la Cour suprême.
\item JPX, connectivity services~; HKEX, hosting services~; SGX, Iris-ST FAQs (juillet 2026)~; ASX, Australian Liquidity Centre~; NSE, co-location
capacity release (2025).
\item CSRC, Administrative Rules for Program Trading in the Securities Market (Trial), 2024, et la consultation de 2023 sur la colocalisation des serveurs de trading
(telle que résumée par JunHe).
\end{omsources}

\section{Exercices}

\begin{exercise}[$\star$]\label{exo:nw:the-asia-pacific-map:1}
Quel est le \omterm{def:nw:the-north-american-map:floor}{plancher de latence} à sens unique entre Singapour et Sydney, dans le vide et dans la fibre~?
\end{exercise}

\begin{exercise}[$\star$]\label{exo:nw:the-asia-pacific-map:2}
En utilisant la \cref{prop:nw:the-asia-pacific-map:seq} avec $c = \qty{10}{\micro\second}$, de combien de temps le dixième membre
du même port reçoit-il un tick après le premier membre~?
\end{exercise}

\begin{exercise}[$\star$]\label{exo:nw:the-asia-pacific-map:3}
Pourquoi l'adresse non publiée du bâtiment d'une plateforme compte-t-elle peu pour Tokyo--Hong Kong, et beaucoup à l'intérieur d'un centre de données~?
\end{exercise}

\begin{exercise}[$\star\star$]\label{exo:nw:the-asia-pacific-map:4}
Un randomiseur tire à nouveau l'ordre des copies à chaque tick. Pourquoi ne réduit-il pas l'écart, et que change-t-il~?
\end{exercise}

\begin{exercise}[$\star\star$]\label{exo:nw:the-asia-pacific-map:5}
Dans le modèle, pourquoi le premier membre d'un tick randomisé se trouve-t-il généralement sur le serveur secondaire~?
\end{exercise}

\begin{exercise}[$\star\star$]\label{exo:nw:the-asia-pacific-map:6}
SGX proposera ITCH sur TCP et en multicast en colocalisation. Lequel choisir pour une stratégie sensible à la latence, et à quoi réserver
l'autre~?
\end{exercise}

\begin{exercise}[$\star\star\star$]\label{exo:nw:the-asia-pacific-map:7}
\emph{Programmation.} Calculer l'avance du premier membre sur le membre médian pour 2, 5, 10 et 20 membres par port, et comparer avec la
\cref{prop:nw:the-asia-pacific-map:seq}.
\end{exercise}

\begin{exercise}[$\star\star\star$]\label{exo:nw:the-asia-pacific-map:8}
\emph{Trouver la faille.} «~Tous les membres colocalisés avaient des câbles de même longueur, donc tous ont reçu le flux tick par tick au même moment.~»
\end{exercise}

\section{Problème~: Premier à se connecter}

\begin{problem}[{Problème du week-end --- la valeur de la connexion en premier}]\label{pb:nw:the-asia-pacific-map:1}
Une plateforme envoie son flux tick par tick en TCP depuis quatre serveurs de diffusion, chacun desservant trois ports de $m$ membres, ainsi qu'un serveur secondaire
peu chargé~; le centre consacre \qty{5}{\micro\second} à chaque copie vers un serveur, et un port \qty{10}{\micro\second} à chaque copie vers un membre. Les membres réagissent
à un tick en \qty{20}{\micro\second} en moyenne (écart type \qty{5}{\micro\second}).

\medskip\noindent\textbf{Partie I --- Un port.}
\begin{enumerate}
\item Avec $m = 10$, quand le premier et le dernier membre du premier port du premier serveur reçoivent-ils un tick~?
\item Combien de temps attend le membre médian de ce port~?
\item Comment cela change-t-il sur le quatrième serveur, dans l'ordre du centre~?
\item Quel membre de l'ensemble de la topologie reçoit en premier, et lequel en dernier~?
\end{enumerate}

\medskip\noindent\textbf{Partie II --- L'avantage.}
\begin{enumerate}[resume]
\item Selon le modèle, quel est l'avantage du premier membre sur le membre médian pour $m = 2, 5, 10, 20$~?
\item Comparer avec la \cref{prop:nw:the-asia-pacific-map:seq}.
\item À quelle fréquence le premier membre devance-t-il le membre médian jusqu'au moteur, tick après tick~?
\item Quelle est l'ampleur de l'avantage face à une course tranchée en quelques microsecondes (livre~11)~?
\end{enumerate}

\medskip\noindent\textbf{Partie III --- Les remèdes.}
\begin{enumerate}[resume]
\item Que fait un randomiseur à la fourchette, à la corrélation de rang et au taux de victoire persistant~?
\item Qu'apporte la fermeture du serveur secondaire au trafic normal~?
\item Que fait le multicast~?
\item Quel remède faudrait-il demander à la plateforme, et pourquoi~?
\end{enumerate}

\medskip\noindent\textbf{Partie IV --- Le verdict.}
\begin{enumerate}[resume]
\item Énoncer le \emph{résultat nommé}~: l'avantage du premier membre sur le membre médian avec dix membres par port, avec et sans multicast.
\item Lesquels des paramètres du modèle sont des faits issus de l'ordonnance de la SEBI, et lesquels sont des hypothèses~?
\item Comment mesurer la fourchette réelle sur le flux d'une plateforme~?
\item Pourquoi la circulaire de 2015 exige-t-elle une latence similaire au regard de l'infrastructure fournie par la bourse, plutôt que des câbles égaux~?
\item Comment la consultation chinoise de 2023 formule-t-elle la même règle~?
\item Où ailleurs la distribution d'une copie par client apparaît-elle dans ce livre~?
\item Quelle est la différence entre un avantage de latence dû à un chemin plus court et un avantage dû à un ordre d'envoi~?
\item En une phrase~: pourquoi la connexion en premier comptait-elle~?
\end{enumerate}
\end{problem}

\section{Questions d'entretien}

\begin{interviewq}[$\star$ \iqroles{développeur}]\label{iq:nw:the-asia-pacific-map:1}
Quelle est la différence entre les données de marché en unicast et en multicast, et pourquoi les bourses préfèrent-elles le multicast en \omterm{def:nw:colocation-products-and-how-they-are-sold:colo}{colocation}~?
\end{interviewq}

\begin{interviewq}[$\star$ \iqroles{développeur, trader}]\label{iq:nw:the-asia-pacific-map:2}
À peu près à quelle distance se trouvent Tokyo, Hong Kong et Singapour en temps de propagation de la lumière, et qu'est-ce que cela implique pour les stratégies inter-marchés~?
\end{interviewq}

\begin{interviewq}[$\star\star$ \iqroles{développeur}]\label{iq:nw:the-asia-pacific-map:3}
Vous recevez le même flux TCP sur deux connexions, et l'une arrive systématiquement 40 microsecondes plus tard. Comment découvririez-vous la cause~?
\end{interviewq}

\begin{interviewq}[$\star\star$ \iqroles{trader, chercheur}]\label{iq:nw:the-asia-pacific-map:4}
Que signifie un accès «~équitable~» aux données de marché, et comment testeriez-vous une plateforme à cet égard~?
\end{interviewq}

\begin{interviewq}[$\star\star$ \iqroles{développeur}]\label{iq:nw:the-asia-pacific-map:5}
Un fournisseur diffuse des cotations à ses clients par websockets depuis un seul serveur. Qui reçoit chaque mise à jour en premier~?
\end{interviewq}

\begin{interviewq}[$\star\star\star$ \iqroles{développeur, chercheur}]\label{iq:nw:the-asia-pacific-map:6}
Concevez une distribution des données de marché pour une nouvelle plateforme qui ne puisse favoriser personne. Que mesurez-vous pour le prouver~?
\end{interviewq}
